\chapter{Crystallographic Background}

\section{Space Groups}

A space group is a mathematical description of the symmetry of a crystal structure. It consists of all symmetry operations (rotations, reflections, inversions, and translations) that leave the crystal invariant.

\subsection{Symmetry Operations}

Each symmetry operation consists of a rotation matrix $\Rmat$ and a fractional translation vector $\mathbf{t}$:

\begin{equation}
    \mathbf{r}' = \Rmat \cdot \mathbf{r} + \mathbf{t}
\end{equation}

where $\mathbf{r}$ is a position vector in direct space.

For 3D space groups, $\Rmat$ is a $3 \times 3$ integer matrix and $\mathbf{t}$ is a 3-vector of fractions. For 2D wallpaper groups, $\Rmat$ is a $2 \times 2$ integer matrix and $\mathbf{t}$ is a 2-vector of fractions (promoted to 3D with z-invariance).

\subsection{Point Group vs. Space Group}

\begin{itemize}
    \item The \textbf{point group} consists only of the rotation/reflection matrices (no translations). It describes the rotational symmetry.
    \item The \textbf{space group} includes both rotations and translations. It describes the full symmetry.
    \item The order of the point group (number of rotation matrices) is always a divisor of the order of the space group (number of symmetry operations).
\end{itemize}

\section{Reciprocal Lattice}

The reciprocal lattice is the Fourier transform of the direct lattice. For a direct lattice with basis vectors $\mathbf{a}_1, \mathbf{a}_2, \mathbf{a}_3$, the reciprocal lattice basis vectors are:

\begin{equation}
    \mathbf{b}_1 = 2\pi \frac{\mathbf{a}_2 \times \mathbf{a}_3}{\mathbf{a}_1 \cdot (\mathbf{a}_2 \times \mathbf{a}_3)}, \quad
    \mathbf{b}_2 = 2\pi \frac{\mathbf{a}_3 \times \mathbf{a}_1}{\mathbf{a}_1 \cdot (\mathbf{a}_2 \times \mathbf{a}_3)}, \quad
    \mathbf{b}_3 = 2\pi \frac{\mathbf{a}_1 \times \mathbf{a}_2}{\mathbf{a}_1 \cdot (\mathbf{a}_2 \times \mathbf{a}_3)}
\end{equation}

Any reciprocal lattice vector can be written as:

\begin{equation}
    \mathbf{q} = h \mathbf{b}_1 + k \mathbf{b}_2 + l \mathbf{b}_3 = (h, k, l)_{\text{recip}}
\end{equation}

where $(h, k, l)$ are the Miller indices (integers).

\subsection{Metric Tensor}

The metric tensor $\Gmat$ encodes the geometry of the reciprocal lattice:

\begin{equation}
    G_{ij} = \mathbf{b}_i \cdot \mathbf{b}_j
\end{equation}

The squared magnitude of a reciprocal lattice vector is:

\begin{equation}
    |\mathbf{q}|^2 = \mathbf{q}^T \Gmat \mathbf{q} = \sum_{i,j} h_i G_{ij} h_j
\end{equation}

\section{Stars}

A \textbf{star} of a reciprocal lattice vector $\mathbf{q}$ is the set of all vectors generated by applying the point group operations:

\begin{equation}
    \text{Star}(\mathbf{q}) = \{ \Rmat \cdot \mathbf{q} \mid \Rmat \in \text{Point Group} \} \cup \{ -\Rmat \cdot \mathbf{q} \mid \Rmat \in \text{Point Group} \}
\end{equation}

The union with negatives ensures the star is closed under inversion, which is required for real-valued Fourier series.

\subsection{Star Properties}

\begin{itemize}
    \item \textbf{Multiplicity}: The number of vectors in the star. Determined by the order of the point group divided by the order of the stabilizer subgroup.
    \item \textbf{Closure}: A star is closed if it contains all symmetry-equivalent vectors. The function \texttt{star\_is\_closed()} verifies this.
    \item \textbf{q-squared}: All vectors in a star have the same $|\mathbf{q}|^2$ (q-squared), since rotations preserve length.
\end{itemize}

\section{Miller Indices}

Miller indices $(h, k, l)$ label reciprocal lattice vectors and crystallographic planes. In the SAFB library, Miller indices are used to:

\begin{itemize}
    \item Identify and group reciprocal lattice vectors into stars
    \item Compute the canonical key for lexicographic sorting
    \item Label modes in the output
\end{itemize}

The canonical key is computed by sorting the Miller indices and taking the absolute value of the first non-zero element, ensuring a unique identifier for each family of planes.
